% ===========================================================================
% Section: Introduction
% ===========================================================================
\section{Introduction}
\label{sec:introduction}

3D Gaussian Splatting (3DGS)~\cite{kerbl2023_3dgs} has emerged as a powerful
primitive for photo-realistic scene representation, achieving real-time
rendering quality that surpasses NeRF-based methods while maintaining
explicit, editable geometry. The integration of 3DGS with simultaneous
localization and mapping (SLAM) pipelines---particularly those fusing
LiDAR, inertial, and camera data---has produced systems such as
\gaussianlic~\cite{lang2026_gaussianlic2} that achieve real-time,
photo-realistic mapping on a single moving platform.

However, existing Gaussian SLAM systems share a critical limitation: they are
\emph{tightly coupled to the Robot Operating System} (ROS) and therefore
restricted to Linux environments. The ROS dependency manifests at multiple
levels: message-passing via \texttt{roscpp}, image transport via
\texttt{cv\_bridge}, point cloud conversion via \texttt{pcl\_conversions},
pose messages via \texttt{geometry\_msgs}, and coordinate transforms via
\texttt{tf}. This coupling prevents deployment on Windows workstations,
limits interoperability with commercial sensor SDKs (e.g., Orbbec,
Intel RealSense), and complicates integration into industrial pipelines
that operate outside the ROS ecosystem.

Furthermore, the real-time design of these systems assumes a continuous
online data stream, making them suboptimal for \emph{offline} batch
processing of pre-recorded datasets---a common workflow in surveying,
inspection, and content production.

We address these limitations with \winlic, a framework that makes the
following contributions:

\begin{enumerate}
    \item \textbf{ROS-free architecture.} We refactor the Gaussian-LIC2
    mapping pipeline to eliminate all ROS dependencies, producing a
    standalone C++ application that builds with pure CMake and runs on
    both Windows (MSVC) and Linux (GCC).

    \item \textbf{Pluggable sensor interface.} We design an abstract
    \texttt{ISensorSource} interface that decouples the reconstruction
    pipeline from the data source. Offline file-based readers and live
    camera SDKs (Orbbec, future sensors) are integrated through a uniform
    four-method interface.

    \item \textbf{Multi-modal depth fusion.} We introduce a
    confidence-weighted depth fusion module that combines sparse LiDAR
    depth, dense iToF depth, and stereo SGM depth into a single
    high-quality depth map, which is then completed by a TensorRT-based
    neural network (SPNet).

    \item \textbf{Preserved algorithm fidelity.} The core Gaussian
    Splatting algorithm---initialization, map extension, optimization,
    and densification---is preserved verbatim from \gaussianlic, ensuring
    reconstruction quality is not compromised by the refactoring.
\end{enumerate}

The remainder of this paper is organized as follows.
Section~\ref{sec:related} reviews related work in Gaussian SLAM and
multi-modal fusion.
Section~\ref{sec:system_overview} presents the system architecture.
Section~\ref{sec:sensor_fusion} details the pluggable sensor interface
and multi-modal depth fusion strategy.
Section~\ref{sec:gaussian_reconstruction} describes the Gaussian
Splatting reconstruction pipeline.
Section~\ref{sec:experiments} discusses the experimental setup, and
Section~\ref{sec:conclusion} concludes.
